\usepackage{booktabs} % Add to preamble if not already there
\usepackage{tabularx} % For flexible column widths

\begin{table}[htbp]
\centering
\caption{Dimensions of informational efficiency for multi-risk assessments. Rather than a single scalar metric, efficiency is operationalised as a multi-dimensional evaluation spanning statistical, decisional, temporal, and institutional domains.}
\label{tab:efficiency_metrics}
\begin{tabularx}{\textwidth}{@{} l l X l @{}}
\toprule
\textbf{Efficiency Dimension} & \textbf{Core Decision Question} & \textbf{Candidate Metric / Proxy} & \textbf{Theoretical Origin} \\
\midrule

\textbf{Information Gain} & Does adding this complexity (e.g., a hazard layer) reduce uncertainty about the risk state? & 
\textit{Mutual Information}: Quantifies the reduction in entropy of the outcome variable given the new model input. \newline
\textit{Variance Reduction}: Percentage decrease in output variance relative to a simpler baseline. & 
Information Theory \newline \newline Statistical Learning \\

\midrule

\textbf{Decision Impact} & Does the added complexity actually change the decision that would be made? & 
\textit{Decision Divergence Rate}: Frequency with which the complex model triggers a different categorical action (e.g., evacuate vs. shelter) compared to a simpler heuristic. \newline
\textit{Expected Value of Including Uncertainty (EVIU)}: Change in expected loss when acting on the complex model vs. a baseline. & 
Decision Analysis \newline \newline Bayesian Decision Theory \\

\midrule

\textbf{Temporal Efficiency} & Does the model yield the decision within the available operational time window? & 
\textit{Decision Latency}: Time elapsed from data input to an actionable, interpretable output. \newline
\textit{Time-to-Threshold}: Ratio of model processing time to the physical lead time of the hazard. & 
Operations Research \newline \newline Early Warning Systems \\

\midrule

\textbf{Institutional Uptake} & Can the end-user actually process and defend the output under pressure? & 
\textit{Utilisation Rate (ITT proxy)}: Proportion of offered model runs that are actually used in the final decision (Intention-to-Treat). \newline
\textit{Cognitive Load}: Number of distinct probabilistic inputs a decision-maker must synthesize to reach a binary action. & 
Behavioural Economics \newline \newline Bounded Rationality \\

\bottomrule
\end{tabularx}
\end{table}